\chapter{Complexity — Counter Prog Poly Time}
%%% generated by the blueprint pipeline from TCSlib.Complexity.ClassNP.CounterProgPolyTime
%%%%%%%%%%%%%%%%%%%%%%%%%%%%%%%%%%%%%%%%%%%%%%%%%%%%%%%%%%%%%%%%%%%%%%%%%%%%%%%
\section{Overview}

This module shows that a counter program (a goto program over finitely many unary
natural-number registers that reads its input once and prints bits) which halts within
polynomially many abstract steps computes a polynomial-time computable function: its
compiled multi-tape Turing machine simulates $t$ abstract steps within $t(2t+3)$ machine
steps. This is how the emitters of [AB09, Remark~6.7] are shown to run in polynomial time.

%%%%%%%%%%%%%%%%%%%%%%%%%%%%%%%%%%%%%%%%%%%%%%%%%%%%%%%%%%%%%%%%%%%%%%%%%%%%%%%
\section{Declarations}

\begin{theorem}[Polynomially many counter-program steps give FP]
\lean{Complexity.CounterProg.polyTimeComputable}
\label{Complexity.CounterProg.polyTimeComputable}
\leanok
\uses{Complexity.CounterProg.exists_tm, Complexity.CounterProg.init, Complexity.CounterProg.run, Complexity.CounterProg.toTM, Complexity.PolyTimeComputable, Turing.FinTM.ComputesInTime.mono}
\difficulty{4}
Let $P$ be a counter program over $R$ natural-number registers with a finite set of labels $\Lambda$, i.e.\ an assignment to each label of one instruction (halt, jump, print a bit, increment a register, decrement a register saturating at $0$, branch on a register being zero, print the value $v$ of a register as $1^v$, or read and consume the next input bit, branching on end-of-input, $0$ or $1$). Fix a start label $\ell_0 \in \Lambda$, a function $f : \{0,1\}^* \to \{0,1\}^*$ and constants $C, c \in \mathbb{N}$, and suppose that for every input $x \in \{0,1\}^*$ there is some $t \le C(|x|+1)^c$ such that running $P$ on $x$ for $t$ steps from the initial state (label $\ell_0$, all registers $0$, no input read, empty output) reaches a halted state whose output is $f(x)$. Then $f$ is polynomial-time computable: some finite binary-alphabet Turing machine computes $f$ within $C'(n+1)^{c'}$ steps on every input of length $n$, for some constants $C', c' \in \mathbb{N}$.
\end{theorem}

\begin{theorem}[Halting counter programs in polynomial time give FP]
\lean{Complexity.CounterProg.polyTimeComputable_of_goes}
\label{Complexity.CounterProg.polyTimeComputable_of_goes}
\leanok
\uses{Complexity.CounterProg.Goes, Complexity.CounterProg.Instr, Complexity.CounterProg.init, Complexity.CounterProg.polyTimeComputable, Complexity.PolyTimeComputable}
\difficulty{3}
Let $P$ be a counter program over $R$ natural-number registers with a finite set of labels $\Lambda$ (each label carrying one instruction: halt, jump, print a bit, increment, saturating decrement, zero-test branch, print a register's value $v$ as $1^v$, or read and consume the next input bit), let $\ell_0 \in \Lambda$, let $f : \{0,1\}^* \to \{0,1\}^*$ and let $C, c \in \mathbb{N}$. Suppose that for every input $x \in \{0,1\}^*$ there are register values $\rho : \{0,\dots,R-1\} \to \mathbb{N}$ and a position $p \in \mathbb{N}$ such that, for every previously printed output $o$, running $P$ on $x$ from label $\ell_0$ with all registers $0$, no input read and output $o$ reaches, within at most $C(|x|+1)^c$ steps, the halted state with registers $\rho$, $p$ input bits read and output $o$ followed by $f(x)$. Then $f$ is polynomial-time computable: some finite binary-alphabet Turing machine computes $f$ within $C'(n+1)^{c'}$ steps on every input of length $n$, for some constants $C', c' \in \mathbb{N}$.
\end{theorem}
